\documentclass[11pt,a4paper]{article}
\usepackage[margin=23mm]{geometry}
\usepackage{amsmath,amssymb,amsthm}
\setlength{\emergencystretch}{2em}
\newtheorem{theorem}{Theorem}
\newtheorem{proposition}[theorem]{Proposition}
\newcommand{\gen}[1]{\langle#1\rangle}
\title{The first later cut when the adjacent suffix is intransitive}
\author{Human Compression project: L044 continuation}
\date{}
\begin{document}\maketitle
Retain the complete critical-head hypotheses and acquired proofs:
$G\le A_n$ is transitive with a connected string C-marking,
all maximal omissions are intransitive, and the actual first three
marks give an extendable six-point head on $D$. In the contrary
range $n=2r\ge18$, write
\[
 P_k=G_{\le k},\qquad R=G_{\ge3}
       =\operatorname{Alt}(\Omega\setminus(D\setminus\{6\})).
\]
The actual head orbit $C_k$ of $P_k$ carries its pure alternating
kernel, with the complete acquired parity-fibre formula and strict
growth. In this note impose the extra hypothesis
\[
                    B=G_{\ge4}\text{ is intransitive on }
                    \Omega\setminus D.                 \tag{1}
\]
Eligibility always refers to a full ambient omission.

\begin{theorem}
Let $5\le j\le r-2$ be the least later label with a unique eligible edge.
Suppose this cut has a first shared label $g<j$ with $g\ge5$,
and $P=P_{g-1}$ is intransitive on its support $Y$.
Then all of the following hold.
\begin{enumerate}
\item $P$ has precisely two nontrivial orbits: its head orbit
      $C_{g-1}$ and one exterior orbit. Moreover $|Y|=2g+1$.
\item $j=g+1$. Write the two $G_j$-orbits as $C=C_{j-1}$
      and $A=\Omega\setminus C$, with crossing endpoint $b\in C$.
      Then
      \[
       |A|=2(r-g-1),\qquad
       \operatorname{Fix}(P)=A\cup\{b\},\qquad
       C=C_{g-1}\mathbin{\dot\cup} E\mathbin{\dot\cup}\{b\},
      \]
      where $E$ is the exterior $P$-orbit.
\item Put $Q=G_{>j}$ and $t=r-g-1$. If $Q$ is transitive
      on $A$, it has degree $2t$ and rank $t-1$ there.
      If it is intransitive on $A$, it has two faithful orbits
      of size $t$, and its action on each is the natural $S_t$.
\item The actual interval $H=G_{\{3,\ldots,g-1\}}$ has
      two nontrivial orbits, $C_{g-1}\setminus(D\setminus\{6\})$
      and $E$. Its paired ambient interval bound is an equality:
      \[
        2(g-3)=|\operatorname{supp}(H)|-2.
      \]
      A minimum-square selection has exactly one tree component
      in each of these two orbits; square components are not excluded.
\end{enumerate}
\end{theorem}
\begin{proof}
The complete early-cut calculation gives at least two eligible
edges at label $3$ under (1), and at label $4$ in every critical
head. Minimality of $j$ gives at least two for every label
$5,\ldots,j-1$: at least one exists because the full group is
transitive and its omission is intransitive. Hence every label
$3,\ldots,g-1$ is paired in the ambient sense.

The acquired unique-cut orbit lemma identifies the head-side
$G_j$-orbit as $C_{j-1}$ and shows that $Q$ fixes it pointwise.
Its other orbit $A$ is nontrivial, because a singleton admits
no shared label. The hypothesis $j\le r-2$ makes the reversed
crossing interior, as required by these support theorems. The full reversed transitive-side or intransitive-side
support theorem therefore applies. It gives
\[
                         |Y|\le2g+1.                    \tag{2}
\]
This is its actual-support conclusion, not a bound on an arbitrary
projected group.

Let $p$ be the number of nontrivial $P$-orbits, so $p\ge2$.
The exact prefix intersection formula, using actual ambient IP,
identifies $H=P\cap R$. It removes exactly the five head points
$D\setminus\{6\}$ from the support and from the head orbit,
and leaves every exterior orbit unchanged. Since $g\ge5$, the
remaining head orbit has at least three points and is nontrivial.
Thus
\[
 |\operatorname{supp}(H)|=|Y|-5,\qquad o^*(H)=p,
\]
where $o^*$ counts nontrivial orbits. The full acquired ambient
interval theorem gives
\[
 2g-6=2|\{3,\ldots,g-1\}|
       \le |Y|-5-p\le2g-4-p.                            \tag{3}
\]
It follows that $p\le2$. Hence $p=2$, and both inequalities
in (3) are equalities. This proves the first and fourth conclusions,
including the tree-component count from equality in that theorem.

Now (2) is an equality. Apply the complete equality interfaces
of the reversed support theorems, with reversed crossing
$i'=r-1-j$ and last shared index $h'=r-1-g$. In the transitive
$Q$ case, they give $h'=i'+1$, $|A|=2h'$, and exact fixed
set $A\cup\{b\}$. In the intransitive $Q$ case, strictness
unless $h'=i'+1$ gives the same index equality; equality forces
two faithful $Q$-orbits of common size $i'+1$, natural symmetric
action on each, and the same exact fixed set. Thus in both cases
$j=g+1$, $h'=t$, and $|A|=2t$.
The literal suffix rank is $r-1-j=t-1$. The exact fixed set
and the two nontrivial $P$-orbits give the asserted partition
of $C$. No conclusion that $Q$ is natural alternating, or that
its transitive case is already classified, has been inserted.
\end{proof}


\begin{proposition}[The bound without adjacent-suffix transitivity information]
Retain the theorem's interior-cut, minimality, $g\ge5$ and
intransitive-prefix hypotheses, but drop (1). Then $P_{g-1}$
has at most three nontrivial orbits. If it has three, then
$|Y|=2g+1$, $j=g+1$, and the same reversed-side equality
conclusions hold. Moreover adjoining $\rho_3$ to
$K=G_{\{4,\ldots,g-1\}}$ merges exactly two of its point orbits
and leaves every other orbit unchanged.
\end{proposition}
\begin{proof}
The label $4$ and all labels $5,\ldots,g-1$ remain paired;
no assertion about label $3$ is needed. Write
$d(V)=n-o(V)$ for the point-orbit deficit of a subgroup $V$.
The acquired ambient interval theorem gives
$d(K)\ge2(g-4)$. Adjoining $\rho_3$ strictly changes the
orbit partition. Otherwise $K$ and $H=\gen{K,\rho_3}$
would have the same point orbits; adjoining all other ambient
marks would then give identical orbit partitions for $G_3$
and $G$, contradicting the intransitive omission. Hence
\[
 d(H)\ge d(K)+1\ge2g-7.
\]
The exact prefix intersection still gives $d(H)=|Y|-5-p$,
where $p$ counts the nontrivial $P$-orbits. With (2) this yields
$p\le3$. At $p=3$ all inequalities are equalities. In particular
$|Y|=2g+1$ and the same reversed interior equality interfaces
apply. The equality $d(H)-d(K)=1$ says precisely that the number
of point orbits decreases by one; as adjoining a generator can
only merge orbits, exactly two old orbits merge.
\end{proof}

\paragraph{Attribution and scope.}
The new step is to combine the complete acquired actual prefix
intersection with the acquired ambient interval deficit before
using the reversed equality interfaces. Their full proofs remain
paid, including all symmetric-string inputs. The underlying source
route is L044 Propositions 5.20--5.32. This is a conditional
reduction for all $g$ at once; it does not eliminate the resulting
configuration. Endpoint cuts $j=r-1$ remain separate: the reversed interior
equality interfaces used here do not apply to them.
The adjacent suffix hypothesis (1) is essential
for the paired label $3$. The case where $P$ is transitive on
its support is governed separately by the reviewed exceptional-end
alternative. The final proposition gives a weaker three-orbit bound without (1).
No whole-original, question, formal or numerical credit
is claimed.
\end{document}
